\section{Discussion}
\label{sec:discuss}

\mypara{Admission control for poor wireless channel conditions.}
While \sys{}'s RAN resource management is independent of wireless channel status, poor channel conditions can still degrade performance. 
LC applications with high bandwidth demands but weak channel quality may consume nearly all wireless resources while missing their SLOs, degrading performance for other users. 
An admission control mechanism can address this by profiling application throughput requirements against UE channel status and terminating service when channel quality is insufficient. 
This preserves SLO satisfaction for UEs with acceptable channel conditions while maintaining efficient spectrum utilization. 
Recent work such as Zipper~\cite{zipper24nsdi} provides techniques for designing such mechanisms.

\mypara{Handling downlink contention.}
Although this work focuses on uplink RAN scheduling, downlink contention is also important to consider. 
Downlink transmissions may suffer from congestion, but the downlink channel is usually more stable and predictable than uplink.
This stability enables end-to-end rate control mechanisms to converge to appropriate sending rates when timely client feedback is available. 
Prior work~\cite{zhuge22sigcomm} leverages uplink ACKs for faster congestion feedback, but delayed uplink grants in 5G networks can hinder timeliness. 
By improving uplink scheduling, \sys{} can prevent excessive queueing at UE buffers that would delay feedback packets (\eg TCP ACKs), ensuring timely congestion signals and enabling more effective downlink rate adaptation.

\mypara{Dealing with UE handover.}
While the applications studied in this work do not involve UE mobility across cells, scenarios such as autonomous driving may trigger inter-cell handovers. 
In deployments where multiple cells are managed under the same base station, handovers can often be absorbed locally without session disruption. 
The more challenging case involves mobility across different base stations, which requires transferring both radio and edge-session state across sites.
To support LC applications under such conditions, we envision proactively replicating UE state across base stations to enable seamless session continuation after handover. 
Designing such mechanisms introduces open questions around synchronization overhead and consistency.

\modified{\mypara{Fairness vs.\ latency trade-off.}
\sys{} prioritizes LC requests to meet their SLOs, consistent with our design goals (\autoref{sec:design-goals}).
Stronger throughput fairness between LC and best-effort (BE) traffic would delay LC transmissions under contention, increasing tail latency and reducing SLO satisfaction.
Our evaluation reflects this tension: systems that prioritize fairness between LC and BE traffic (\eg \tutti{} and \arma{}) incur more LC SLO violations than \sys{}, as shown in \autoref{fig:network-latency-steady-workload} and \autoref{fig:network-latency-dynamic-workload}.
Accordingly, for BE workloads \sys{} provides starvation freedom rather than strong throughput fairness under heavy contention.
Achieving strong fairness guarantees while preserving tight latency bounds for LC applications under wireless contention remains an open problem.


\mypara{Limitations.}
Our request identification approach at the RAN is constrained by the aggregation semantics of BSR.
When requests arrive in rapid succession within a single BSR reporting interval, they appear as a single aggregated increase at the MAC layer.
In such cases, \sys{} can only infer a shared request start time and performs scheduling at the request-group level rather than achieving per-request granularity.
This limitation affects scenarios with extremely high request rates but does not impact typical workloads where inter-request intervals exceed BSR reporting periods.

A second limitation concerns processing time estimation for compute resource scheduling.
\sys{} relies on processing time estimates to proactively allocate compute resources, which works well for applications with stable processing pipelines like those in our evaluation.
However, applications with dynamic processing patterns (\eg adaptive VR rendering) can exhibit high variance in processing latency, making accurate estimation challenging.
In such cases, estimation errors may reduce the effectiveness of \sys{}'s compute-side scheduling, though RAN-side scheduling remains unaffected and continues to minimize network delays.}